% year: תש"פ
% type: מבחן
% semester: א
% moed: ב
% part: ב
% number: 1
% points: 9
% categories: derivatives, continuity, multiple-choice
% source: exams/מבחנים/תשפ/מועד ב תשפ.pdf

\begin{question}
תהי $f(x)$ פונקציה כך ש-$f(x)+\sin(x)$ גזירה בקטע הפתוח $(a,b)$. מה מהבאים חייב להיות נכון?
\begin{enumerate}
  \item $f(x)+\cos(x)$ גזירה בקטע הפתוח $(a,b)$.
  \item $f(x)+\sin(x)$ רציפה בקטע הסגור $[a,b]$.
  \item $f(x)$ רציפה בקטע הסגור $[a,b]$.
  \item $f(x)$ חסומה בקטע הפתוח $(a,b)$.
  \item $f(x)+\sin(x)$ חסומה בקטע הפתוח $(a,b)$.
\end{enumerate}
\end{question}

\begin{hint}
כתבו את $f+\cos$ כסכום של $f+\sin$ ופונקציה גזירה ידועה.
\end{hint}

\begin{solution}
\textbf{התשובה הנכונה: א.}

נסמן $g(x)=f(x)+\sin x$, גזירה ב-$(a,b)$. אז
\[ f(x)+\cos x=g(x)-\sin x+\cos x , \]
וזה סכום של פונקציות גזירות ב-$(a,b)$ ($\sin$ ו-$\cos$ גזירות בכל $\R$), ולכן גזיר ב-$(a,b)$. כלומר א' נכונה תמיד.

שאר הטענות אינן חייבות להתקיים. דוגמה נגדית: $f(x)=\frac{1}{x-a}-\sin x$ ב-$(a,b)$. אז $f(x)+\sin x=\frac{1}{x-a}$ גזירה ב-$(a,b)$, אבל:
\begin{itemize}
  \item $f+\sin$ ו-$f$ אינן מוגדרות (ובוודאי אינן רציפות) ב-$x=a$ --- ב' וג' נופלות. (גם אם נגדיר להן ערך כלשהו ב-$a$, הן אינן חסומות ליד $a$ ולכן אינן רציפות בה.)
  \item $\frac{1}{x-a}\to+\infty$ כאשר $x\to a^+$, ולכן $f+\sin$ אינה חסומה ב-$(a,b)$ --- ה' נופלת; וגם $f=\frac{1}{x-a}-\sin x$ אינה חסומה (כי $\sin$ חסומה) --- ד' נופלת.
\end{itemize}
\end{solution}
